\chapter{\texorpdfstring{$\beta$}{Beta}-Elimination}\label{ch:b1:elimination}

Treat 2-bromo-2-methylpropane with sodium ethoxide in ethanol at room
temperature and the main product is an ether; heat the same mixture and
the main product is an alkene, 2-methylpropene, with no ethoxy group in
it. The ethoxide ion has acted not as a \omterm{def:b1:electronic-effects:nucleophile}{nucleophile}, attacking carbon, but
as a base, removing a proton from the carbon next door while the bromide
left. This competition between substitution and elimination runs through
all of organic synthesis. This chapter describes the two elimination
mechanisms, their remarkable geometric requirement, the rule that
predicts which alkene forms, and how to steer a reaction one way or the
other.

\begin{recall}
SN1 and SN2 (\cref{ch:b1:nucleophilic-substitution}); \omterm{def:b1:stereochemistry-in-depth:representations}{Newman projections},
\omterm{def:b1:stereochemistry-in-depth:isomers}{conformations} of cyclohexane and $Z$/$E$ descriptors
(\cref{ch:b1:stereochemistry-in-depth}); \omterm{def:b1:electronic-effects:intermediates}{carbocations} and bases
(\cref{ch:b1:electronic-effects}).
\end{recall}

\section{The E2 mechanism}

\begin{definition}[\texorpdfstring{$\beta$}{Beta}-elimination, E2 and E1]\label{def:b1:elimination:elimination}
A \emph{$\beta$-elimination}\index{$\beta$-elimination} removes a
\omterm{def:b1:electronic-effects:nucleophile}{leaving group} Y from a carbon (the $\alpha$ carbon) and a proton from a
neighbouring carbon (a $\beta$ carbon), forming a double bond between
them. In the \emph{E2 mechanism}\index{E2 mechanism} a base takes the
proton while Y leaves, in one \omterm{def:b1:elementary-steps:elementary-step}{elementary step}; in the \emph{E1
mechanism}\index{E1 mechanism} Y first leaves, giving a \omterm{def:b1:electronic-effects:intermediates}{carbocation}, which
then loses a $\beta$ proton.
\end{definition}

\begin{proposition}[Rate law of E2]\label{prop:b1:elimination:e2-law}
For E2, $v = k[\mathrm{RY}][\mathrm{B}]$, where B is the base.
\end{proposition}

\begin{proof}
One bimolecular \omterm{def:b1:elementary-steps:elementary-step}{elementary step}, involving \omterm{def:b1:nucleophilic-substitution:substitution}{substrate} and base: the law of
mass action of \omterm{def:b1:elementary-steps:elementary-step}{elementary steps} (\cref{ch:b1:elementary-steps}).
\end{proof}

\begin{definition}[Anti- and syn-periplanar]\label{def:b1:elimination:periplanar}
Two bonds on neighbouring atoms are \emph{anti-periplanar}\index{anti-periplanar}
when they lie in one plane on opposite sides (\omterm{def:b1:stereochemistry-in-depth:conformations}{torsion angle} $180^\circ$),
\emph{syn-periplanar}\index{syn-periplanar} when they lie in one plane on
the same side ($0^\circ$).
\end{definition}

\begin{omfigure}
\begin{tikzpicture}[font=\small, >={Stealth}]
  \pic (n) at (0,0) {omnewman={90}{270}};
  \node at (n-f1) {H}; \node at (n-f2) {\ce{CH3}}; \node at (n-f3) {H};
  \node at (n-b1) {Br}; \node at (n-b2) {\ce{CH3}}; \node at (n-b3) {H};
  \node (b) at (-2.6,2.0) {\ce{EtO-}};
  \draw[->, thick, omThm] (b) to[bend left=20] ($(n-f1)+(-0.15,0.1)$);
  \draw[->, thick, omThm] (0.12,0.8) to[bend left=40] (0.12,0.25);
  \draw[->, thick, omThm] (0.12,-0.6) to[bend right=35] ($(n-b1)+(0.25,0.15)$);
  \node[below] at (0,-2.0) {H and Br anti-periplanar};
  \draw[->, thick] (2.4,0) -- (3.6,0) node[midway, above] {E2};
  \node at (5.6,0) {\chemfig{C(-[:150]H_3C)(-[:210]H)=C(-[:30]H)(-[:-30]CH_3)}};
  \node at (5.6,-0.9) {\cip{E}-but-2-ene};
\end{tikzpicture}

{\small E2 reaction of 2-bromobutane, viewed along the C3--C2 bond (C3 in
front, C2 bearing Br behind) in one of its two \omterm{def:b1:elimination:periplanar}{anti-periplanar}
\omterm{def:b1:stereochemistry-in-depth:isomers}{conformations}; the other one, with the other hydrogen of C3 anti to Br,
gives \cip{Z}-but-2-ene. Three electron
pairs move together: the base takes the proton, the C--H pair becomes the
$\pi$ bond, the C--Br pair leaves with bromine. The groups that face each
other in the \omterm{def:b1:stereochemistry-in-depth:representations}{Newman projection} end on the same side of the double bond.}
\end{omfigure}

\begin{proposition}[E2 is anti and stereospecific]\label{prop:b1:elimination:anti}
E2 proceeds from the \omterm{def:b1:stereochemistry-in-depth:isomers}{conformation} in which the C--H and C--Y bonds are
\omterm{def:b1:elimination:periplanar}{anti-periplanar}. It is stereospecific: diastereomeric \omterm{def:b1:nucleophilic-substitution:substitution}{substrates} give
different alkene \omterm{def:b1:stereochemistry-in-depth:isomers}{stereoisomers}, fixed by the anti arrangement.
\end{proposition}

\begin{proof}
In the \omterm{def:b1:elementary-steps:energy-profile}{transition state} the C--H \omterm{def:b1:lewis-resonance-vsepr:lewis}{bonding pair} must enter the
antibonding region of the C--Y bond while the two carbons become
trigonal; the overlap is best when the two bonds are parallel, in one
plane. The anti arrangement is staggered (low energy) and keeps the base
far from the \omterm{def:b1:electronic-effects:nucleophile}{leaving group}; the syn one is eclipsed. With H and Y anti,
the two other groups on each carbon are fixed: those that are on the same
side in the \omterm{def:b1:stereochemistry-in-depth:representations}{Newman projection} end \textit{cis} on the double bond. A
diastereomer has two groups exchanged on one carbon, and gives the other
alkene.
\end{proof}

\begin{proposition}[E2 on a cyclohexane]\label{prop:b1:elimination:cyclohexane-e2}
On a cyclohexane ring, E2 requires the \omterm{def:b1:electronic-effects:nucleophile}{leaving group} and the $\beta$
hydrogen to be both \omterm{def:b1:stereochemistry-in-depth:conformations}{axial} (\textit{trans}-diaxial); a \omterm{def:b1:electronic-effects:nucleophile}{leaving group} held
\omterm{def:b1:stereochemistry-in-depth:conformations}{equatorial} cannot react until the \omterm{def:b1:stereochemistry-in-depth:conformations}{ring flips}.
\end{proposition}

\begin{proof}
On adjacent carbons of a \omterm{def:b1:stereochemistry-in-depth:conformations}{chair}, two \omterm{def:b1:stereochemistry-in-depth:conformations}{axial} bonds are \omterm{def:b1:elimination:periplanar}{anti-periplanar} (one
up, one down, parallel to the axis); an \omterm{def:b1:stereochemistry-in-depth:conformations}{equatorial} bond is anti to ring
C--C bonds only, never to a C--H bond.
\end{proof}

\begin{omfigure}
\begin{tikzpicture}[font=\small]
  \pic (a) at (0,0) {omchair};
  \draw[thick, omThm] (a-c1) -- (a-a1) node[above] {Cl};
  \draw[thick] (a-c1) -- (a-e1) node[right] {H};
  \draw[thick, omDef] (a-c2) -- (a-a2) node[below] {H};
  \draw[thick, omDef] (a-c6) -- (a-a6) node[below] {H};
  \node[align=left, anchor=west] at (2.6,0.2) {Cl axial (up) on C1;\\axial H (down) on both neighbours:\\two anti-periplanar H};
\end{tikzpicture}

{\small \textit{trans}-Diaxial arrangement on a \omterm{def:b1:stereochemistry-in-depth:conformations}{chair}: the \omterm{def:b1:stereochemistry-in-depth:conformations}{axial} chlorine
is \omterm{def:b1:elimination:periplanar}{anti-periplanar} to the \omterm{def:b1:stereochemistry-in-depth:conformations}{axial} hydrogens of the two neighbouring carbons
(blue).}
\end{omfigure}

\section{The E1 mechanism}

\begin{proposition}[Rate law of E1]\label{prop:b1:elimination:e1-law}
For E1, $v = k[\mathrm{RY}]$, the rate being that of the ionisation.
\end{proposition}

\begin{proof}
As for SN1 (\cref{prop:b1:nucleophilic-substitution:sn1-law}): the
\omterm{def:b1:electronic-effects:intermediates}{carbocation} is formed in a slow step and consumed fast, here by loss of a
proton to the solvent or a \omterm{def:b1:predominant-reaction:strong-acid}{weak base}; the \omterm{def:b1:elementary-steps:qssa}{steady-state approximation}
leaves $v = k_1[\mathrm{RY}]$.
\end{proof}

E1 shares its \omterm{def:b1:electronic-effects:intermediates}{carbocation} with SN1: tertiary \omterm{def:b1:nucleophilic-substitution:substitution}{substrates} in \omterm{def:b1:intermolecular-forces-solvents:solvent-classes}{protic
solvents} give both, the alkene being favoured by heating. Since the
\omterm{def:b1:electronic-effects:intermediates}{carbocation} is planar and the proton is lost afterwards, E1 has no
geometric requirement and is not stereospecific; it gives mainly the more
stable alkene, usually the \cip{E} one.

\begin{omfigure}
\setchemfig{atom sep=2.1em}
\schemestart
\chemfig{C(-[2]CH_3)(-[4]H_3C)(-[6]CH_3)-Br}
\arrow{->[slow]}
\chemfig{C^{+}(-[2]CH_3)(-[:210]H_3C)(-[:-30]CH_3)}
\arrow{->[fast][\ce{-H+}]}
\chemfig{C(-[2]CH_3)(-[:210]H_3C)=[:-30]CH_2}
\schemestop
\par\vspace{4pt}

{\small E1: ionisation of 2-bromo-2-methylpropane to the \omterm{def:b1:electronic-effects:intermediates}{carbocation}, then
loss of a proton from a methyl group to the solvent, giving
2-methylpropene.}
\end{omfigure}

\section{Regioselectivity: Zaitsev's rule}

\begin{definition}[Regioselective and stereoselective reactions]\label{def:b1:elimination:selectivity}
A reaction that can give several \omterm{def:b1:stereochemistry-in-depth:isomers}{constitutional isomers} but forms one
mainly is a \emph{regioselective reaction}\index{regioselective reaction};
one that can give several \omterm{def:b1:stereochemistry-in-depth:isomers}{stereoisomers} but forms one mainly is a
\emph{stereoselective reaction}\index{stereoselective reaction}.
\end{definition}

\begin{proposition}[Zaitsev's rule]\label{prop:b1:elimination:zaitsev}
When several $\beta$ hydrogens can be removed, the major alkene is
usually the most substituted (the most stable), and the \cip{E} isomer
rather than the \cip{Z} (\emph{Zaitsev's rule}\index{Zaitsev's rule}).
Bulky bases, and geometric constraints such as the trans-diaxial
requirement, can override it.
\end{proposition}

\begin{proof}
Alkyl groups on the carbons of a double bond stabilise it (as they
stabilise \omterm{def:b1:electronic-effects:intermediates}{carbocations}), and \cip{E} alkenes avoid the crowding of
\textit{cis} groups. The \omterm{def:b1:elementary-steps:energy-profile}{transition states} of E2 and E1 have partial
double-bond character, so the more stable alkene is formed faster. A
bulky base reaches the less hindered hydrogens more easily; and a
hydrogen that cannot be \omterm{def:b1:elimination:periplanar}{anti-periplanar} to the \omterm{def:b1:electronic-effects:nucleophile}{leaving group} is not
removed by E2 at all, whatever the stability of the alkene it would give.
\end{proof}

\begin{method}[Predicting the alkene]\label{met:b1:elimination:predict-alkene}
\begin{enumerate}
  \item List the $\beta$ carbons that carry hydrogens.
  \item For E2, keep only hydrogens that can be \omterm{def:b1:elimination:periplanar}{anti-periplanar} to Y (on a
        ring: trans-diaxial, in a \omterm{def:b1:stereochemistry-in-depth:conformations}{chair} that can actually form).
  \item Among them, favour the one that gives the most substituted alkene
        (Zaitsev), unless the base is bulky.
  \item For each, draw the \omterm{def:b1:stereochemistry-in-depth:representations}{Newman projection} in the \omterm{def:b1:stereochemistry-in-depth:conformations}{anti conformation} to
        read the \cip{E}/\cip{Z} geometry.
\end{enumerate}
\end{method}

\section{Substitution or elimination?}

\begin{proposition}[Substitution against elimination]\label{prop:b1:elimination:sn-vs-e}
Elimination is favoured by strong, bulky bases, by high temperature and by
tertiary or crowded \omterm{def:b1:nucleophilic-substitution:substitution}{substrates}; substitution by good \omterm{def:b1:electronic-effects:nucleophile}{nucleophiles} that are
\omterm{def:b1:predominant-reaction:strong-acid}{weak bases}, by low temperature and by primary \omterm{def:b1:nucleophilic-substitution:substitution}{substrates}.
\end{proposition}

\begin{proof}
Base strength favours the attack on hydrogen; bulk hinders the attack on
a crowded carbon but not on an exposed hydrogen. A \omterm{def:b1:electronic-effects:carbon-class}{tertiary carbon} is
inaccessible to SN2, and its \omterm{def:b1:electronic-effects:intermediates}{carbocation} loses protons easily. Elimination
turns one molecule into two or three (alkene, \omterm{def:b1:electronic-effects:nucleophile}{leaving group}, protonated
base): it is favoured as the temperature rises, a result justified in
the Year~2 volume.
\end{proof}

\begin{omfigure}
\begin{tabular}{llll}
\hline
\omterm{def:b1:nucleophilic-substitution:substitution}{substrate} & \omterm{def:b1:predominant-reaction:strong-acid}{strong base}, small & \omterm{def:b1:predominant-reaction:strong-acid}{strong base}, bulky & \omterm{def:b1:predominant-reaction:strong-acid}{weak base}, good \omterm{def:b1:electronic-effects:nucleophile}{nucleophile} \\
\hline
primary & SN2 (some E2) & E2 & SN2 \\
secondary & E2 (and SN2) & E2 & SN2 (aprotic), SN1/E1 (protic) \\
tertiary & E2 & E2 & SN1/E1 (protic) \\
\hline
\end{tabular}

\smallskip
{\small A first orientation: the main reaction in each case. Heating
favours the elimination column by column.}
\end{omfigure}

\section{Exercises}

\begin{exercise}[$\star$]\label{exo:b1:elimination:1}
Give the alkenes that can form by elimination of HBr from 2-bromobutane,
from 2-bromo-2-methylbutane and from bromocyclohexane, and the one \omterm{prop:b1:elimination:zaitsev}{Zaitsev's
rule} predicts as major.
\end{exercise}

\begin{exercise}[$\star$]\label{exo:b1:elimination:2}
Write the \omterm{def:b1:rate-laws:order}{rate laws} of E1 and E2. How would you tell them apart
experimentally for 2-bromo-2-methylpropane in ethanol containing ethoxide?
\end{exercise}

\begin{exercise}[$\star$]\label{exo:b1:elimination:3}
Draw the \omterm{def:b1:stereochemistry-in-depth:representations}{Newman projections} of 2-bromobutane along C2--C3 in its three
\omterm{def:b1:stereochemistry-in-depth:conformations}{staggered conformations}, and say which can undergo E2.
\end{exercise}

\begin{exercise}[$\star$]\label{exo:b1:elimination:4}
Predict the main reaction (SN1, SN2, E1, E2): bromoethane with sodium
hydroxide in water at \qty{25}{\celsius}; 2-bromo-2-methylpropane with
potassium tert-butoxide; 2-bromopropane with sodium iodide in propanone;
2-bromo-2-methylpropane in hot ethanol.
\end{exercise}

\begin{exercise}[$\star\star$]\label{exo:b1:elimination:5}
Show, with \omterm{def:b1:stereochemistry-in-depth:representations}{Newman projections}, that the E2 elimination of HBr from the
two \omterm{def:b1:stereochemistry-in-depth:enantiomers}{diastereomers} of 2-bromo-3-methylpentane that can lose the C3
hydrogen gives the two \omterm{def:b1:stereochemistry-in-depth:isomers}{stereoisomers} of 3-methylpent-2-ene.
\end{exercise}

\begin{exercise}[$\star\star$]\label{exo:b1:elimination:6}
Write the mechanism of the dehydration of 2-methylbutan-2-ol in hot
concentrated sulfuric acid (E1) and give the major alkene.
\end{exercise}

\begin{exercise}[$\star\star$]\label{exo:b1:elimination:7}
Explain why \textit{trans}-1-bromo-4-tert-butylcyclohexane reacts very
slowly by E2 while its \textit{cis} isomer reacts fast. (The tert-butyl
group stays \omterm{def:b1:stereochemistry-in-depth:conformations}{equatorial}.)
\end{exercise}

\begin{exercise}[$\star\star$]\label{exo:b1:elimination:8}
2-Bromo-2-methylbutane with sodium ethoxide gives mainly
2-methylbut-2-ene; with potassium tert-butoxide, mainly
2-methylbut-1-ene. Explain.
\end{exercise}

\begin{exercise}[$\star\star$]\label{exo:b1:elimination:9}
Why is elimination never observed with bromomethane? With
1-bromo-2,2-dimethylpropane by E2?
\end{exercise}

\begin{exercise}[$\star\star\star$]\label{exo:b1:elimination:10}
In ethanol at \qty{25}{\celsius}, 2-bromo-2-methylpropane gives a mixture
of the ether (SN1) and the alkene (E1) whose proportions do not depend on
the concentration of the \omterm{def:b1:nucleophilic-substitution:substitution}{substrate}. Show that both products come from
the same \omterm{def:b1:electronic-effects:intermediates}{carbocation} and that their ratio is the ratio of two \omterm{def:b1:rate-laws:order}{rate
constants}.
\end{exercise}

\begin{exercise}[$\star\star\star$]\label{exo:b1:elimination:11}
One diastereomer of 1,2-dibromo-1,2-diphenylethane (the meso one) gives,
by E2 with ethoxide, a single stereoisomer of 1-bromo-1,2-diphenylethene.
Find which, with a \omterm{def:b1:stereochemistry-in-depth:representations}{Newman projection}.
\end{exercise}

\begin{exercise}[$\star\star\star$]\label{exo:b1:elimination:12}
Using the energies of \cref{ch:b1:stereochemistry-in-depth}, explain why
an E2 reaction on a cyclohexane can be slowed by a factor of hundreds
when the \omterm{def:b1:electronic-effects:nucleophile}{leaving group} must become \omterm{def:b1:stereochemistry-in-depth:conformations}{axial} in a crowded \omterm{def:b1:stereochemistry-in-depth:conformations}{chair}. Express the
rate as the product of a conformer fraction and a \omterm{def:b1:rate-laws:order}{rate constant}.
\end{exercise}

\section{Problem: Menthyl and Neomenthyl Chlorides}

\begin{problem}[{Weekend problem --- the chairs of two diastereomeric
chlorides, the anti-periplanar hydrogens available in each, the alkenes
formed, and the ratio of their E2 rates}]
\label{pb:b1:elimination:1}
Menthyl chloride and neomenthyl chloride are the chlorides of menthol
and neomenthol (\cref{ch:b1:stereochemistry-in-depth}): on the ring, C1
carries Cl, C2 the isopropyl group, C5 the methyl group. In menthyl
chloride the three groups can all be \omterm{def:b1:stereochemistry-in-depth:conformations}{equatorial}; in neomenthyl chloride,
Cl is \omterm{def:b1:stereochemistry-in-depth:conformations}{axial} when the two alkyl groups are \omterm{def:b1:stereochemistry-in-depth:conformations}{equatorial}. Both are treated
with sodium ethoxide in ethanol, an E2 reaction. Data of the problem:
\qty{95}{\percent} of neomenthyl chloride is in the \omterm{def:b1:stereochemistry-in-depth:conformations}{chair} with Cl \omterm{def:b1:stereochemistry-in-depth:conformations}{axial};
for menthyl chloride, the \omterm{def:b1:stereochemistry-in-depth:conformations}{chair} with Cl \omterm{def:b1:stereochemistry-in-depth:conformations}{axial} (all three groups \omterm{def:b1:stereochemistry-in-depth:conformations}{axial})
is a fraction \num{5.0e-3}; assume that each \omterm{def:b1:elimination:periplanar}{anti-periplanar} hydrogen of a
reactive \omterm{def:b1:stereochemistry-in-depth:conformations}{chair} reacts with the same \omterm{def:b1:rate-laws:order}{rate constant}. Neomenthyl chloride
gives \qty{78}{\percent} of the trisubstituted alkene.

\medskip
\noindent\textbf{Part I --- The \omterm{def:b1:stereochemistry-in-depth:conformations}{chairs}.}
\begin{enumerate}
  \item Draw the preferred \omterm{def:b1:stereochemistry-in-depth:conformations}{chair} of menthyl chloride.
  \item Draw the preferred \omterm{def:b1:stereochemistry-in-depth:conformations}{chair} of neomenthyl chloride.
  \item Are the two chlorides \omterm{def:b1:stereochemistry-in-depth:enantiomers}{enantiomers} or \omterm{def:b1:stereochemistry-in-depth:enantiomers}{diastereomers}?
  \item State the geometric requirement of E2 on a ring.
  \item In which \omterm{def:b1:stereochemistry-in-depth:conformations}{chair} can menthyl chloride react? Draw it.
  \item Why is that \omterm{def:b1:stereochemistry-in-depth:conformations}{chair} rare? Use the energy cost of \omterm{def:b1:stereochemistry-in-depth:conformations}{axial} groups.
\end{enumerate}

\medskip
\noindent\textbf{Part II --- The anti hydrogens.}
\begin{enumerate}[resume]
  \item List the $\beta$ carbons of the chloride (C2 and C6) and their
        hydrogens.
  \item In the reactive \omterm{def:b1:stereochemistry-in-depth:conformations}{chair} of neomenthyl chloride, which hydrogens are
        \omterm{def:b1:elimination:periplanar}{anti-periplanar} to Cl?
  \item In the reactive \omterm{def:b1:stereochemistry-in-depth:conformations}{chair} of menthyl chloride, is the hydrogen of C2
        \omterm{def:b1:stereochemistry-in-depth:conformations}{axial}? Is a hydrogen of C6 \omterm{def:b1:stereochemistry-in-depth:conformations}{axial}?
  \item How many \omterm{def:b1:elimination:periplanar}{anti-periplanar} hydrogens does each chloride offer?
  \item Draw a \omterm{def:b1:stereochemistry-in-depth:representations}{Newman projection} along C1--C6 for the reactive \omterm{def:b1:stereochemistry-in-depth:conformations}{chair} of
        menthyl chloride.
  \item Why could a \omterm{def:b1:elimination:periplanar}{syn-periplanar} elimination not rescue the
        reaction?
\end{enumerate}

\medskip
\noindent\textbf{Part III --- The products.}
\begin{enumerate}[resume]
  \item Name the alkene formed by removing the C2 hydrogen and the one
        formed by removing a C6 hydrogen.
  \item Which one does \omterm{prop:b1:elimination:zaitsev}{Zaitsev's rule} favour, and why?
  \item What does neomenthyl chloride give? Is the 78/22 ratio consistent
        with \omterm{prop:b1:elimination:zaitsev}{Zaitsev's rule}?
  \item What does menthyl chloride give? Is this consistent with
        \omterm{prop:b1:elimination:zaitsev}{Zaitsev's rule}?
  \item Which reaction is stereospecific, regioselective, or both?
  \item Why is substitution (SN2) minor with ethoxide here?
\end{enumerate}

\medskip
\noindent\textbf{Part IV --- The rates.}
\begin{enumerate}[resume]
  \item Write the rate of each reaction as a function of the fraction of
        reactive \omterm{def:b1:stereochemistry-in-depth:conformations}{chair} and the number of \omterm{def:b1:elimination:periplanar}{anti-periplanar} hydrogens.
  \item Compute the relative rate of neomenthyl chloride.
  \item Compute the relative rate of menthyl chloride.
  \item Compute the ratio of the rates $k(\text{neomenthyl})/k(\text{menthyl})$.
\end{enumerate}
\end{problem}
